\documentclass[11pt]{article}
\usepackage[margin=1in]{geometry}
\usepackage{amsmath, amssymb, amsthm}
\usepackage{hyperref}
\usepackage{enumitem}
\usepackage{graphicx}
\usepackage{bm}
\usepackage{mathrsfs}
\usepackage{listings}
\usepackage{courier}

\lstset{
  basicstyle=\ttfamily\small,
  breaklines=true,
  columns=fullflexible
}

\title{Multiplicity Wire Harness for Neuroplasticity:\\
Prime-Indexed Cognitive Routing with Golden-Ratio Entropy Bound}
\author{Citizen Gardens \\ The Foundation of Multiplicity}
\date{April 27, 2026}

\begin{document}
\maketitle

\begin{abstract}
This defensive publication discloses a formal, implementable architecture---the \emph{Multiplicity Wire Harness}---for modeling neuroplasticity and cognition as a prime-indexed, recursively evolving state space with built-in ethical constraints. Building on an initial neuroplasticity framework in Multiplicity Theory \cite{neuroplasticityMT}, we reinterpret primes not as cognitive primitives but as \emph{addresses} in a sparse cognitive space, define a recursive update operator $\Xi$ over this address space, introduce a Consciousness Stability Law (CSL) that bounds changes in address-entropy by $\Delta S < \ln \varphi$ (where $\varphi$ is the golden ratio), and specify a neuroquantum interface mapping internal states to EEG-like observables via a Hermitian operator $\hat{H}_\Xi$ and inverse Fourier transform. We further formalize an ASD-centric braiding construction (EchoBraid), propose a golden-ratio--motivated stability interpretation, and give an RFC-style harness specification (MT--HARNESS--001) separating harness from payload. Example pseudocode illustrates practical implementation using sparse prime-indexed dictionaries and entropy enforcement. This document is intended to establish prior art for the conceptual, mathematical, and implementation-level components of this architecture.
\end{abstract}
\newpage
\tableofcontents
\newpage
\section{Executive Summary}

\subsection{Purpose and Scope}

This publication discloses a general-purpose \emph{cognitive routing architecture}---the \textbf{Multiplicity Wire Harness}---that:

\begin{itemize}[noitemsep]
  \item Models cognitive/neuroplastic states as elements of a prime-indexed tensor product Hilbert space.
  \item Uses a recursive update operator $\Xi$ to model learning as lawful, stepwise evolution rather than ad hoc reconfiguration.
  \item Embeds an \emph{ethical safety constraint}---the Consciousness Stability Law (CSL)---as a bound on address entropy $\Delta S < \ln \varphi$, inspired by golden-ratio organization and entropy considerations in EEG. \cite{neuroplasticityMT,EEGEntropyForce,EEGEntropyBiomarker,EEGEntropyASD}
  \item Provides a neuroquantum interface: a mapping from internal states to EEG-like signals via a Hermitian operator $\hat{H}_\Xi$ and inverse Fourier transform.
  \item Includes an ASD-aligned \emph{EchoBraid} construction for neurodivergent routing modes.
\end{itemize}

Crucially, this architecture is specified as a \emph{harness}, not a specific cognitive theory. Any model that can be expressed in terms of prime-indexed traces $\theta_p$, attention phases $\phi_p$, and a compatible $\Xi$ can be \emph{plugged into} this harness.

\subsection{Key Contributions}

\begin{enumerate}[noitemsep]
  \item \textbf{Prime-Indexed Cognitive Address Space}: We interpret primes as \emph{addresses} indexing channels in a Hilbert space, not as semantic atoms. This yields a unique, sparse addressing scheme for cognitive traces, leveraging the uniqueness of prime factorization.

  \item \textbf{Recursive Cognitive Dynamics}: Learning is modeled as a recursive update $\Theta(t+1) = \Xi(\Theta(t))$, where $\Theta(t)$ encodes prime-indexed content and phases, and $\Xi$ preserves the address structure while updating amplitudes and phases.

  \item \textbf{Consciousness Stability Law (CSL)}: We impose a bound on address-entropy change $\Delta S < \ln \varphi$, where $\varphi$ is the golden ratio, linking psychological safety and trauma-awareness to an information-theoretic constraint on the dispersion of activation over addresses. \cite{neuroplasticityMT,EEGEntropyForce,GoldenEEG,EEGEntropyBiomarker}

  \item \textbf{Neuroquantum Interface}: We define an observable learning signal (EEG-like) as
  \[
  \mathrm{EEG}_{\mathrm{Learn}}(t) = \mathcal{F}^{-1}\big(\langle \Psi(t) \vert \hat{H}_\Xi \vert \Psi(t) \rangle\big),
  \]
  linking the recursive cognitive state to spectral brain signals via a Hermitian operator derived from $\Xi$. \cite{neuroplasticityMT,QuantumEEG1,QuantumEEG2}

  \item \textbf{EchoBraid for ASD}: For ASD-aligned cognition, we introduce EchoBraid, a braided superposition of identity states over ordered primes, and we connect this to empirical findings on EEG entropy and complexity in autism. \cite{neuroplasticityMT,EEGEntropyASD,EEGEntropyASD2}

  \item \textbf{MT--HARNESS--001 Specification}: We present an RFC-style specification with MUST/MAY/SHOULD requirements, separating harness (addressing, recursion, CSL, interface) from payload (semantics of $\theta_p$, specific $\Xi$, concrete $\hat{H}_\Xi$).
\end{enumerate}

\subsection{Prior Art Context}

Existing work has explored:

\begin{itemize}[noitemsep]
  \item Entropy of EEG signals as markers of cortical excitability, force modulation, seizure susceptibility, emotion regulation, and ASD-related alterations. \cite{EEGEntropyForce,EEGEntropySeizure,EEGEntropyBiomarker,EEGEntropyASD,EEGEntropyASD2}
  \item Golden ratio organization in EEG, suggesting $\varphi$-based frequency ratios may support optimal decoupling and desynchronization. \cite{GoldenEEG}
  \item Entropy-based control of cognitive models and neural systems. \cite{EntropyCognitiveModels,EntropyNeuralControl}
\end{itemize}

However, to our knowledge, no prior work has:

\begin{itemize}[noitemsep]
  \item Treated primes as a \emph{computational addressing scheme} for cognitive traces in a recursive Hilbert space model.
  \item Embedded a golden-ratio--based entropy bound as an intrinsic \emph{ethical constraint} on learning trajectories.
  \item Specified an implementable harness layer (MT--HARNESS--001) that decouples routing/safety from theoretical payloads.
\end{itemize}

This document therefore establishes a new, combinatorial-mathematical substrate for cognitive modeling and neuroplasticity, with clear paths to empirical validation.

\section{Mathematical Overview}

\subsection{Prime-Indexed Cognitive Space}

Let $\mathbb{P}$ denote the set of prime numbers:
\[
\mathbb{P} = \{2,3,5,7,11,\dots\}.
\]

Each prime $p \in \mathbb{P}$ is interpreted as a \emph{cognitive address}. We associate to each $p$ a two-dimensional complex Hilbert space $\mathcal{H}_p \cong \mathbb{C}^2$. The full cognitive Hilbert space is given by the (possibly infinite) tensor product:
\[
\mathcal{H} = \bigotimes_{p \in \mathbb{P}} \mathcal{H}_p.
\]

In practice, at any finite time $t$, only a finite subset $\mathbb{P}_{\mathrm{active}}(t) \subset \mathbb{P}$ is active. Implementations therefore use a sparse representation over this active set.

\subsection{State Encoding: Traces and Phases}

The cognitive state at time $t$ is encoded as
\begin{equation}
\Psi(t) = \sum_{p \in \mathbb{P}_{\mathrm{active}}(t)} \theta_p(t) \otimes e^{i \phi_p(t)},
\end{equation}
where:

\begin{itemize}[noitemsep]
  \item $\theta_p(t) \in \mathcal{H}_p$ is the \emph{cognitive trace} at address $p$ (content/amplitude).
  \item $\phi_p(t) \in \mathbb{R}$ is the \emph{attention phase} at address $p$.
\end{itemize}

This structure was introduced in the original neuroplasticity document within Multiplicity Theory, under the heading of Prime-Indexed Recursive Tensor Mathematics (PIRTM). \cite{neuroplasticityMT}

\subsection{Recursive Operator and Learning Dynamics}

Let $\Theta(t)$ denote the collection of all address-specific states at time $t$:
\[
\Theta(t) = \{\theta_p(t), \phi_p(t)\}_{p \in \mathbb{P}_{\mathrm{active}}(t)}.
\]

Learning is modeled as a self-updating recursive function:
\begin{equation}
\Theta(t+1) = \Xi(\Theta(t)),
\end{equation}
where $\Xi$ is a (possibly nonlinear) operator. \cite{neuroplasticityMT}

We require:

\begin{enumerate}[label=(\alph*),noitemsep]
  \item \textbf{Address preservation}: If a trace is indexed by $p$ at time $t$, its successor at $t+1$ is also indexed by $p$ (the harness preserves addresses).
  \item \textbf{Deterministic evolution}: $\Xi$ is deterministic given $\Theta(t)$.
  \item \textbf{Locality with coupling}: $\Xi$ may couple multiple addresses, but such coupling must be explicitly defined (e.g., via arithmetic relations among primes).
\end{enumerate}

A multi-scale extension introduces $\{\Xi_n\}_{n=1}^N$ acting at different temporal resolutions:
\[
\Theta_n(t+1) = \Xi_n(\Theta_n(t), \Theta_{n-1}(t)),
\]
with $\Theta_0(t)$ representing a base level (e.g., fast synaptic processes).

\subsection{Address Entropy}

We define an \emph{address entropy} over the active prime set at time $t$.

Let
\[
w_p(t) = \|\theta_p(t)\|^2 \quad\text{for }p \in \mathbb{P}_{\mathrm{active}}(t),
\]
and
\[
Z(t) = \sum_{p \in \mathbb{P}_{\mathrm{active}}(t)} w_p(t).
\]
Assuming $Z(t) > 0$, define normalized address weights:
\[
\pi_p(t) = \frac{w_p(t)}{Z(t)}.
\]

Then the address entropy is:
\begin{equation}
S[\Psi(t)] = -\sum_{p \in \mathbb{P}_{\mathrm{active}}(t)} \pi_p(t) \ln \pi_p(t).
\end{equation}

This is the Shannon entropy of the probability distribution over active addresses, analogous to entropy used in cognitive and neural models. \cite{EntropyCognitiveModels,EntropyNeuralControl}

\subsection{Consciousness Stability Law (CSL)}

The Consciousness Stability Law (CSL) imposes a bound on the change in address entropy during a learning episode.

For a learning episode $[t, t+\Delta t]$, define:
\[
\Delta S = S[\Psi(t+\Delta t)] - S[\Psi(t)].
\]

The CSL requires:
\begin{equation}
\Delta S < \ln \varphi,
\end{equation}
where $\varphi = \frac{1 + \sqrt{5}}{2}$ is the golden ratio. \cite{neuroplasticityMT}

Interpretation:

\begin{itemize}[noitemsep]
  \item $\Delta S$ measures how much the distribution of activation over addresses spreads or contracts.
  \item Bounding $\Delta S$ prevents sudden, unbounded dispersion of activation across many new addresses, which is hypothesized to correlate with psychological overload or trauma.
  \item The use of $\ln \varphi$ relates to proposals that golden-ratio organization provides optimal decoupling or desynchronization in neural systems and may partition known/unknown information in a scale-invariant manner. \cite{GoldenEEG,GoldenRatioInfo}
\end{itemize}

Later sections discuss prior-art support for golden-ratio organization in EEG and entropy modulation. \cite{GoldenEEG,EEGEntropyForce,EEGEntropyBiomarker}

\subsection{Neuroquantum Interface}

We define a (neuro)quantum-like interface mapping cognitive state evolution to observable EEG-like signals. In the original document, neuroplastic changes were modeled as frequency-domain shifts:
\[
\mathrm{EEG}_{\mathrm{Learn}}(t) = \mathcal{F}^{-1}(\langle \psi \vert \hat{H} \vert \phi \rangle),
\]
where $\hat{H}$ is a Hermitian operator and $\mathcal{F}^{-1}$ denotes inverse Fourier transform. \cite{neuroplasticityMT}

In our refined formulation, we set:
\begin{equation}
\mathrm{EEG}_{\mathrm{Learn}}(t) = \mathcal{F}^{-1}\left(\big\langle \Psi(t) \big| \hat{H}_\Xi \big| \Psi(t) \big\rangle\right),
\end{equation}
where $\hat{H}_\Xi$ is a Hermitian operator constructed from:

\begin{itemize}[noitemsep]
  \item The recursive operator $\Xi$.
  \item Prime-specific potentials $V_p(\theta_p, \phi_p)$.
  \item Possible cross-address couplings.
\end{itemize}

This establishes a bidirectional mapping:

\begin{itemize}[noitemsep]
  \item From $\Psi(t)$ and $\Xi$ to EEG-like time series via $\hat{H}_\Xi$ and $\mathcal{F}^{-1}$.
  \item From EEG data back to constraints on $\Psi(t)$ and $\Xi$, enabling inverse modeling.
\end{itemize}

Quantum-inspired EEG models and entropy analyses already suggest that spectral and entropic features of EEG track changes in excitability, force modulation, seizures, and emotion regulation, making this interface empirically supportable. \cite{QuantumEEG1,QuantumEEG2,EEGEntropyForce,EEGEntropySeizure,EEGEntropyBiomarker,EEGEntropyEmotion}

\subsection{EchoBraid for ASD and Neurodiversity}

The initial neuroplasticity document introduced EchoBraid as an ASD-centric construction:
\begin{equation}
\mathrm{EchoBraid}(t) = \sum_{n=1}^{\infty} \psi_{p_n}(t) \otimes e^{i \theta_{p_n}(t)},
\end{equation}
where $\{p_n\}$ is an ordered sequence of primes. \cite{neuroplasticityMT}

In this defensive publication, we generalize EchoBraid via a neurodiversity parameter $\nu \in [0,1]$:
\begin{equation}
\mathrm{EchoBraid}_\nu(t) = \sum_{n=1}^{\infty} \psi_{p_n}^{1-\nu}(t) \otimes e^{i \theta_{p_n}^{\nu}(t)}.
\end{equation}

Here:

\begin{itemize}[noitemsep]
  \item $\nu = 0$ corresponds to a ``neurotypical'' routing style.
  \item $\nu \to 1$ corresponds to ASD-aligned routing (or other neurodivergent modes).
  \item For all $\nu$, CSL remains in effect, ensuring $\Delta S < \ln \varphi$ for learning episodes.
\end{itemize}

Empirically, EEG studies have shown characteristic differences in entropy, complexity, and spectral features in ASD vs typical development, suggesting that such a parameterization could be detectable. \cite{EEGEntropyASD,EEGEntropyASD2,EEGEntropyASD3}

\section{MT--HARNESS--001: Prime-Indexed Cognitive Routing Layer}

We now present the harness specification as an RFC-style document. This is the core of the \emph{Multiplicity Wire Harness}.

\subsection{Abstract}

MT--HARNESS--001 defines a prime-indexed routing layer for cognitive states with recursive update, an address entropy stability constraint, and a mapping to EEG-like observables. The harness is payload-agnostic: any theory that can be expressed in terms of prime-indexed content $\theta_p$, phases $\phi_p$, a recursive operator $\Xi$, and an appropriate Hermitian operator $\hat{H}_\Xi$ can conform.

\subsection{Definitions and Notation (Harness-Level)}

\begin{itemize}[noitemsep]
  \item $\mathbb{P}$: set of primes (address space).
  \item $\mathbb{P}_{\mathrm{active}}(t)$: finite set of active primes at time $t$.
  \item $\mathcal{H}_p \cong \mathbb{C}^2$: Hilbert space at address $p$.
  \item $\mathcal{H} = \bigotimes_{p \in \mathbb{P}} \mathcal{H}_p$: total cognitive space.
  \item $\theta_p(t)$: content at address $p$ and time $t$.
  \item $\phi_p(t)$: attention phase at address $p$ and time $t$.
  \item $\Psi(t) = \sum_{p \in \mathbb{P}_{\mathrm{active}}} \theta_p(t) \otimes e^{i \phi_p(t)}$: cognitive state at time $t$.
  \item $\Theta(t)$: full configuration of $\{\theta_p(t), \phi_p(t)\}$.
  \item $\Xi$: recursive update operator, $\Theta(t+1) = \Xi(\Theta(t))$.
  \item $S[\Psi(t)]$: address entropy at time $t$.
  \item $\varphi$: golden ratio, $\varphi = \frac{1 + \sqrt{5}}{2}$.
  \item $\hat{H}_\Xi$: Hermitian operator derived from $\Xi$ and the harness.
  \item $\mathrm{EEG}_{\mathrm{Learn}}(t)$: observable learning signal (EEG-like).
\end{itemize}

\subsection{Required Interfaces (MUST)}

\paragraph{Prime addressability.}

\begin{enumerate}[label=(M\arabic*),noitemsep]
  \item Each active cognitive trace MUST be assigned a unique prime address $p \in \mathbb{P}$. No two distinct active traces may share the same prime.
  \item Implementations MUST realize $\mathcal{H}$ via a sparse representation over $\mathbb{P}_{\mathrm{active}}(t)$.
\end{enumerate}

\paragraph{Recursive update.}

\begin{enumerate}[label=(M\arabic*),resume,noitemsep]
  \item Implementations MUST implement a recursive update
  \[
  \Theta(t+1) = \Xi(\Theta(t)),
  \]
  mapping valid harness states to valid harness states.
  \item The operator $\Xi$ MUST preserve prime addresses: traces indexed by $p$ at time $t$ remain indexed by $p$ at time $t+1$.
  \item Any cross-address coupling induced by $\Xi$ MUST be explicitly specified.
\end{enumerate}

\paragraph{CSL safety constraint.}

\begin{enumerate}[label=(M\arabic*),resume,noitemsep]
  \item For any learning episode $[t, t+\Delta t]$, implementations MUST compute $\Delta S = S[\Psi(t+\Delta t)] - S[\Psi(t)]$ using the normalized active set $\mathbb{P}_{\mathrm{active}}(t)$.
  \item Implementations MUST enforce $\Delta S < \ln \varphi$ either by:
  \begin{itemize}[noitemsep]
    \item rejecting or rolling back transitions that violate the bound, or
    \item adjusting the update (e.g., scaling $\Delta \theta_p$) so that the final $\Delta S$ respects the bound.
  \end{itemize}
  \item Implementations MUST log CSL-induced modifications for auditing.
\end{enumerate}

\paragraph{Measurement mapping.}

\begin{enumerate}[label=(M\arabic*),resume,noitemsep]
  \item Implementations MUST provide a mapping
  \[
  \mathrm{EEG}_{\mathrm{Learn}}(t) = \mathcal{F}^{-1}\left(\langle \Psi(t) \vert \hat{H}_\Xi \vert \Psi(t) \rangle\right).
  \]
  \item The operator $\hat{H}_\Xi$ MUST be explicitly defined as a functional of $\Xi$ and the prime addresses (e.g., via prime-local potentials and couplings).
  \item The resulting $\mathrm{EEG}_{\mathrm{Learn}}(t)$ MUST be real-valued, time-indexed, and band-limited in a way that is comparable to empirical EEG data.
\end{enumerate}

\subsection{Optional Extensions (MAY)}

\paragraph{Multi-scale recursion.}

Implementations MAY define a hierarchy $\{\Xi_n\}_{n=1}^N$:
\[
\Theta_n(t+1) = \Xi_n(\Theta_n(t), \Theta_{n-1}(t)),
\]
for modeling different timescales of neuroplasticity.

\paragraph{Adaptive phase coherence.}

Implementations MAY introduce a phase-coherence operator $\mathcal{C}_p(t)$ and learning functional $\mathcal{L}(\Psi)$:
\[
\phi_p(t+1) = \phi_p(t) + \alpha \cdot \mathcal{C}_p(t) \cdot \nabla_{\phi_p} \mathcal{L}(\Psi(t)).
\]

\paragraph{Neurodiversity spectrum.}

Implementations MAY define a neurodiversity parameter $\nu \in [0,1]$ and EchoBraid$_\nu(t)$:
\[
\mathrm{EchoBraid}_\nu(t) = \sum_{n=1}^{\infty} \psi_{p_n}^{1-\nu}(t) \otimes e^{i \theta_{p_n}^\nu(t)},
\]
with CSL enforced for all $\nu$.

\paragraph{Prime-to-frequency pinout.}

Implementations MAY adopt a default pinout mapping:
\[
f_p = p \cdot f_0,
\]
with $f_0 \approx 2.5\,\text{Hz}$, so that $p=2$ maps near theta, $p=3$ near low alpha, $p=5$ near upper alpha/sensorimotor, etc., aligning with known EEG bands. \cite{GoldenEEG,EEGEntropyASD2}

Alternative mappings (e.g., logarithmic) are allowed, but MUST be documented.

\subsection{Validation Tests (SHOULD)}

Implementations SHOULD attempt the following tests:

\begin{itemize}[noitemsep]
  \item \textbf{Prime-frequency learning correlation}: Test whether learning performance correlates with changes in prime-indexed spectral bands above null baselines. \cite{EEGEntropyForce}
  \item \textbf{CSL critical behavior}: Test for optimality of learning near $\Delta S \approx \ln \varphi$, and investigate phase coherence scaling behavior around this bound. \cite{GoldenEEG,GoldenRatioInfo}
  \item \textbf{Neurodiversity signatures}: For EchoBraid$_\nu$, test whether EEG connectivity and entropy measures differ systematically with $\nu$ in ways consistent with ASD vs typical development findings. \cite{EEGEntropyASD,EEGEntropyASD2,EEGEntropyASD3}
\end{itemize}

\section{CSL and Golden-Ratio Justification (Conceptual Sketch)}

A full derivation is beyond this initial defensive publication, but we outline the key motivations.

\subsection{Entropy and Cognitive Stability}

Entropy-based measures have been used to analyze and control cognitive models, neural systems, and EEG dynamics. \cite{EntropyCognitiveModels,EntropyNeuralControl,EEGEntropyForce,EEGEntropyEmotion} In EEG, entropy modulation tracks changes in excitability, force control, seizure likelihood, and emotional regulation. \cite{EEGEntropyForce,EEGEntropySeizure,EEGEntropyBiomarker,EEGEntropyEmotion}

We therefore regard address entropy $S[\Psi]$ as a proxy for:

\begin{itemize}[noitemsep]
  \item The dispersion of cognitive activation across the address space.
  \item The complexity and variability of the current cognitive configuration.
\end{itemize}

Bounding $\Delta S$ per learning episode prevents abrupt, high-amplitude reconfiguration across many addresses at once.

\subsection{Golden Ratio as Stability Partition}

Empirical work suggests that golden-ratio frequency relationships in EEG correspond to desynchronized, resting configurations and may minimize pathological cross-frequency locking. \cite{GoldenEEG} Information-theoretic work has also highlighted golden-ratio partitions $p : (1-p) = 1 : p$ as special self-similar splits between known and unknown information. \cite{GoldenRatioInfo}

Interpreting $\Delta S$ as a shift in the partition between ``known'' (high-weight addresses) and ``unknown'' (low-weight or inactive addresses), a bound of $\ln \varphi$ can be seen as selecting a regime where:

\begin{itemize}[noitemsep]
  \item The system remains within a stable, self-similar balance between consolidation and novelty.
  \item Learning episodes do not induce catastrophic readdressing that would correspond to trauma or destabilization.
\end{itemize}

We therefore choose $\ln \varphi$ as a conservative, structurally motivated threshold on address entropy change.

\section{Code Snippet: Minimal Harness Implementation}

This section provides a pseudocode-style implementation sketch in Python-like syntax. It is not optimized, but demonstrates how one might realize the harness.

\subsection{State Representation}

We represent the state as a sparse dictionary mapping primes to $(\theta_p, \phi_p)$.

\begin{lstlisting}[language=Python,caption={Prime-indexed harness state representation}]
import math
import cmath
import numpy as np

# State: mapping from prime p -> (theta_p, phi_p)
# theta_p: complex amplitude (or small vector)
# phi_p: float phase in radians

State = dict  # type alias

def active_primes(state: State):
    return list(state.keys())
\end{lstlisting}

\subsection{Entropy Computation and CSL Enforcement}

\begin{lstlisting}[language=Python,caption={Address entropy and CSL enforcement}]
GOLDEN_RATIO = (1 + 5**0.5) / 2
LN_PHI = math.log(GOLDEN_RATIO)

def compute_entropy(state: State) -> float:
    # Compute |theta_p|^2
    weights = [abs(theta_p)**2 for (theta_p, phi_p) in state.values()]
    Z = sum(weights)
    if Z <= 0:
        return 0.0
    probs = [w / Z for w in weights]
    # Shannon entropy over active addresses
    return -sum(p * math.log(p) for p in probs if p > 0)

def interpolate_states(state_t: State, state_t1: State, alpha: float) -> State:
    """
    Simple linear interpolation in amplitude space as a placeholder
    for entropy-constrained step scaling.
    """
    new_state = {}
    for p in state_t.keys():
        theta0, phi0 = state_t[p]
        theta1, phi1 = state_t1.get(p, (0+0j, phi0))
        theta_new = theta0 + alpha * (theta1 - theta0)
        phi_new = phi0 + alpha * (phi1 - phi0)
        new_state[p] = (theta_new, phi_new)
    return new_state

def csl_enforce(state_t: State, candidate_state: State,
                ln_phi: float = LN_PHI) -> (State, bool):
    S_t = compute_entropy(state_t)
    S_candidate = compute_entropy(candidate_state)
    delta_S = S_candidate - S_t
    if delta_S < ln_phi:
        return candidate_state, False  # no modification
    
    # Simple backtracking line search on step size alpha in (0,1]
    alpha = 1.0
    for _ in range(20):
        alpha *= 0.5
        adjusted = interpolate_states(state_t, candidate_state, alpha)
        S_adj = compute_entropy(adjusted)
        if (S_adj - S_t) < ln_phi:
            return adjusted, True
    # Fallback: return original state_t if no safe step found
    return state_t, True
\end{lstlisting}

This implementation uses a simple line-search-like scheme to scale the update step until $\Delta S < \ln \varphi$. More sophisticated trust-region methods can be substituted.

\subsection{Recursive Update Skeleton}

\begin{lstlisting}[language=Python,caption={Recursive update skeleton}]
def Xi(state_t: State) -> State:
    """
    Payload-dependent recursive operator.
    Must preserve prime keys and update (theta_p, phi_p).
    For demonstration, we apply a simple rotation and small mixing.
    """
    new_state = {}
    # Example: global small phase increment
    dphi = 0.01
    for p, (theta_p, phi_p) in state_t.items():
        # Example amplitude update: small nonlinearity
        theta_new = theta_p * (1.0 + 0.01j)  # toy dynamics
        phi_new = phi_p + dphi
        new_state[p] = (theta_new, phi_new)
    return new_state

def step_with_csl(state_t: State) -> (State, bool):
    candidate = Xi(state_t)
    adjusted, modified = csl_enforce(state_t, candidate)
    return adjusted, modified
\end{lstlisting}

\subsection{Synthetic EEG Generation}

\begin{lstlisting}[language=Python,caption={Synthetic EEG from harness state}]
def H_operator(p: int, theta_p: complex, phi_p: float, f: float) -> complex:
    """
    Example Hermitian contribution at frequency f from address p.
    Payload-specific; here, a simple phase-coded term.
    """
    # Example: weight by amplitude and a frequency-dependent factor
    return theta_p * cmath.exp(1j * (phi_p - 2*math.pi*f)) / (1 + p)

def eeg_from_state(state: State, freqs: np.ndarray) -> np.ndarray:
    """
    Construct a frequency-domain representation from state and
    apply inverse FFT to obtain EEG-like signal.
    """
    spectrum = []
    for f in freqs:
        val = 0+0j
        for p, (theta_p, phi_p) in state.items():
            val += H_operator(p, theta_p, phi_p, f)
        spectrum.append(val)
    spectrum = np.array(spectrum, dtype=complex)
    eeg_signal = np.fft.ifft(spectrum).real
    return eeg_signal
\end{lstlisting}

This code demonstrates how the harness can be connected to an EEG-like observable via a Hermitian-like operator, in line with the neuroquantum interface concept. \cite{neuroplasticityMT,QuantumEEG1,QuantumEEG2}

\section{Suggested Extensions and Future Work}

\subsection{CSL Derivation and Theorems}

Future work can formalize CSL via:

\begin{itemize}[noitemsep]
  \item Proving existence and uniqueness of $\Xi$ under CSL constraints.
  \item Deriving bounds on learning rates in terms of $\Delta S$ and $\ln \varphi$.
  \item Connecting CSL to stability regions in nonlinear dynamical systems with golden-ratio-like spectral organization. \cite{StabilityRegions}
\end{itemize}

\subsection{Empirical Validation Pathways}

We recommend:

\begin{enumerate}[noitemsep]
  \item Synthetic simulations exploring learning stability as a function of $\Delta S$ and CSL enforcement.
  \item Retrospective EEG analysis of learning tasks using prime-indexed frequency bands and entropy measures. \cite{EEGEntropyForce,EEGEntropyBiomarker}
  \item ASD vs typical comparisons using EchoBraid-inspired connectivity and entropy features. \cite{EEGEntropyASD,EEGEntropyASD2,EEGEntropyASD3}
  \item Testing golden-ratio-based coupling and decoupling hypotheses in resting vs active states. \cite{GoldenEEG}
\end{enumerate}

\section{Conclusion}

This defensive publication has disclosed:

\begin{itemize}[noitemsep]
  \item The conceptual reframing of a Multiplicity-based neuroplasticity framework as a \emph{Multiplicity Wire Harness}.
  \item A mathematical formalization of prime-indexed address spaces, recursive cognitive dynamics, address entropy, CSL, and a neuroquantum EEG interface.
  \item An RFC-style harness specification (MT--HARNESS--001) suitable for direct implementation in computational models.
  \item Example code sketches demonstrating how to realize the harness in practice.
\end{itemize}

By publishing these structures, we establish prior art on the combination of prime-indexed cognitive addressing, golden-ratio--bounded entropy changes, and EEG-linked recursive dynamics as a unified architecture for neuroplasticity and cognition.

\bibliographystyle{plain}
\begin{thebibliography}{10}

\bibitem{neuroplasticityMT}
Citizen Gardens Research Network.
\newblock \emph{Neuroplasticity in Multiplicity Theory: A Formal Overview}, 2025.
\newblock (Thread attachment: Neuroplasticity.pdf).

\bibitem{EEGEntropyForce}
R.~Author et~al.
\newblock Entropy in electroencephalographic signals modulates with force magnitude.
\newblock \emph{Journal/PMC}, 2024. \url{https://pmc.ncbi.nlm.nih.gov/articles/PMC11449659/}.

\bibitem{EEGEntropySeizure}
X.~Author et~al.
\newblock Entropy-based seizure detection from EEG.
\newblock \emph{Frontiers in Neurology}, 2023. \url{https://www.frontiersin.org/articles/10.3389/fnhum.2022.1084061/full}.

\bibitem{EEGEntropyBiomarker}
Y.~Author et~al.
\newblock Entropy in scalp EEG as a biomarker for intervention response.
\newblock \emph{Scientific Reports}, 2023. \url{https://www.nature.com/articles/s41598-023-46113-z}.

\bibitem{EEGEntropyEmotion}
Z.~Author et~al.
\newblock EEG entropy modulation as a biomarker of emotion regulation and cognitive flexibility.
\newblock \emph{Review}, 2025. \url{https://pmc.ncbi.nlm.nih.gov/articles/PMC12639470/}.

\bibitem{EEGEntropyASD}
A.~Author et~al.
\newblock EEG entropy analysis in autistic children.
\newblock \emph{PubMed}, 2019. \url{https://pubmed.ncbi.nlm.nih.gov/30503641/}.

\bibitem{EEGEntropyASD2}
B.~Author et~al.
\newblock Entropy, complexity, and spectral features of EEG signals in autism.
\newblock \emph{Neuropsychiatric Journal}, 2025. \url{https://pmc.ncbi.nlm.nih.gov/articles/PMC11832502/}.

\bibitem{EEGEntropyASD3}
C.~Author et~al.
\newblock An encoder-ensemble fusion model for autism diagnosis using multi-domain EEG signal analysis.
\newblock \emph{GitHub: EEG\_Autism}, 2024. \url{https://github.com/abdul-rasool/EEG_Autism}.

\bibitem{QuantumEEG1}
D.~Author et~al.
\newblock Quantum theory of EEG with application to single-trial ERP.
\newblock \emph{Journal of Neurobiology}, 2021. \url{https://www.lidsen.com/journals/neurobiology/neurobiology-05-01-084}.

\bibitem{QuantumEEG2}
E.~Author et~al.
\newblock Quantum physics analysis of EEG signals for drug effects.
\newblock \emph{Neuroimage}, 2020. \url{https://www.sciencedirect.com/science/article/pii/S0149763414001778}.

\bibitem{GoldenEEG}
F.~Author et~al.
\newblock Golden ratio organization in human EEG is associated with theta-alpha coupling.
\newblock \emph{Human Brain Mapping}, 2026. \url{https://pmc.ncbi.nlm.nih.gov/articles/PMC12996120/}.

\bibitem{EntropyCognitiveModels}
G.~Author.
\newblock The use of entropy for analysis and control of cognitive models.
\newblock \emph{Technical Report}, 2003. \url{http://act-r.psy.cmu.edu/wordpress/wp-content/uploads/2012/12/415belavkinR03.pdf}.

\bibitem{EntropyNeuralControl}
H.~Author et~al.
\newblock Control of entropy in neural models of environmental state.
\newblock \emph{eLife}, 2019. \url{https://elifesciences.org/articles/39404}.

\bibitem{GoldenRatioInfo}
I.~Author et~al.
\newblock A golden-ratio partition of information and the balance between prediction and surprise.
\newblock \emph{arXiv preprint}, 2026. \url{https://arxiv.org/html/2602.15266v1}.

\bibitem{StabilityRegions}
J.~Author.
\newblock Stability regions of nonlinear dynamical systems.
\newblock \emph{Cambridge University Press}, 2015. \url{https://hal.science/hal-02984348v1/file/Regions_of_attraction.pdf}.

\end{thebibliography}

\section*{Mathematical Appendix}

\subsection*{Purpose}

This appendix records the explicit mathematical statements, proof sketches, and operator-norm bounds implied by the harness formalism developed in the article. It is written at the level of a defensive publication: each result is stated clearly enough to establish priority and to guide later formal refinement. The appendix builds directly on the prime-indexed recursive tensor model, the Consciousness Stability Law (CSL), and the Hermitian EEG interface introduced in the main text and source document. \cite{neuroplasticityMT,EntropyCognitiveModels,EntropyNeuralControl,QuantumEEG1,GoldenEEG}

\subsection*{A.1. Basic Setup}

Let $\mathbb{P}_{\mathrm{active}}(t)$ be the finite active prime set at time $t$, and define the finite-dimensional Hilbert space
\[
\mathcal{H}_t \;=\; \bigotimes_{p \in \mathbb{P}_{\mathrm{active}}(t)} \mathcal{H}_p,
\qquad
\mathcal{H}_p \cong \mathbb{C}^2.
\]

A cognitive state is represented by
\[
\Psi(t) \;=\; \sum_{p \in \mathbb{P}_{\mathrm{active}}(t)} \theta_p(t)\otimes e^{i\phi_p(t)},
\]
with amplitudes $\theta_p(t)\in \mathcal{H}_p$ and phases $\phi_p(t)\in\mathbb{R}$. \cite{neuroplasticityMT}

Define weights
\[
w_p(t) := \|\theta_p(t)\|^2,
\qquad
Z(t) := \sum_{p\in\mathbb{P}_{\mathrm{active}}(t)} w_p(t),
\qquad
\pi_p(t) := \frac{w_p(t)}{Z(t)}.
\]
Assume $Z(t)>0$. Then the address entropy is
\[
S[\Psi(t)] := -\sum_{p\in\mathbb{P}_{\mathrm{active}}(t)} \pi_p(t)\ln \pi_p(t).
\]

The recursive update is
\[
\Theta(t+1)=\Xi(\Theta(t)),
\]
where $\Theta(t)$ denotes the full collection of amplitudes and phases. \cite{neuroplasticityMT}

\subsection*{A.2. Entropy Well-Posedness}

\textbf{Proposition A.1 (Well-posedness of address entropy).}
If $\mathbb{P}_{\mathrm{active}}(t)$ is finite and $Z(t)>0$, then $S[\Psi(t)]$ is well-defined, finite, and satisfies
\[
0 \le S[\Psi(t)] \le \ln |\mathbb{P}_{\mathrm{active}}(t)|.
\]

\textbf{Proof.}
Because $\mathbb{P}_{\mathrm{active}}(t)$ is finite and $\pi_p(t)\ge 0$ with
\[
\sum_{p\in\mathbb{P}_{\mathrm{active}}(t)} \pi_p(t)=1,
\]
the collection $\{\pi_p(t)\}$ is a probability distribution on a finite set. The Shannon entropy of any finite probability distribution is finite and nonnegative. Its maximum occurs for the uniform distribution $\pi_p = 1/N$ on $N=|\mathbb{P}_{\mathrm{active}}(t)|$ points, yielding
\[
S_{\max}=-\sum_{p=1}^N \frac{1}{N}\ln\frac{1}{N} = \ln N.
\]
The minimum $0$ occurs when all mass is concentrated on a single address. \hfill $\square$

\subsection*{A.3. Lipschitz Bound for Entropy Under Amplitude Perturbations}

We now bound the change in address entropy under small changes in the normalized address distribution.

\textbf{Proposition A.2 (Entropy continuity bound).}
Let $\pi=(\pi_p)$ and $\tilde{\pi}=(\tilde{\pi}_p)$ be two probability distributions on the same finite active prime set of size $N$. Then
\[
|S(\pi)-S(\tilde{\pi})|
\le
\|\pi-\tilde{\pi}\|_1 \,\ln\!\Big(\frac{N}{\|\pi-\tilde{\pi}\|_1}\Big)
\]
for sufficiently small $\|\pi-\tilde{\pi}\|_1$, and in particular there exists a constant $C_N>0$ such that
\[
|S(\pi)-S(\tilde{\pi})|
\le
C_N \|\pi-\tilde{\pi}\|_1.
\]

\textbf{Proof sketch.}
This is a standard continuity property of Shannon entropy on a finite alphabet. Because the entropy map is continuous on the simplex and locally Lipschitz away from the boundary, there exists a finite constant $C_N$ depending only on the dimension. The more explicit logarithmic modulus follows from classical entropy continuity estimates on finite sets. \hfill $\square$

\textbf{Corollary A.2.1.}
If the recursive update $\Xi$ induces
\[
\|\pi(t+1)-\pi(t)\|_1 \le \delta_t,
\]
then
\[
|S[\Psi(t+1)]-S[\Psi(t)]|
\le
C_N \delta_t.
\]
Hence any step-size control on the normalized weights induces direct control on $\Delta S$.

This justifies trust-region-style CSL enforcement in implementations. \cite{EntropyNeuralControl,StabilityRegions}

\subsection*{A.4. CSL Feasibility Under Step Scaling}

Let $\Theta^\star(t+1)$ be a candidate next state produced by the unconstrained recursive operator, and define an interpolated family
\[
\Theta_\alpha(t+1) := \Theta(t) + \alpha\big(\Theta^\star(t+1)-\Theta(t)\big),
\qquad \alpha\in[0,1].
\]

\textbf{Theorem A.3 (Existence of a CSL-admissible scaled step).}
Assume the map
\[
\alpha \mapsto S[\Psi_\alpha(t+1)]
\]
is continuous on $[0,1]$. Then for every candidate step there exists $\alpha_0>0$ such that for all $\alpha\in[0,\alpha_0)$,
\[
S[\Psi_\alpha(t+1)] - S[\Psi(t)] < \ln \varphi.
\]

\textbf{Proof.}
At $\alpha=0$ we have $\Theta_\alpha(t+1)=\Theta(t)$, so
\[
S[\Psi_0(t+1)] - S[\Psi(t)] = 0.
\]
Since $0<\ln\varphi$, continuity implies there exists $\alpha_0>0$ such that
\[
|S[\Psi_\alpha(t+1)]-S[\Psi(t)]| < \ln\varphi
\]
for all sufficiently small $\alpha$. In particular,
\[
S[\Psi_\alpha(t+1)] - S[\Psi(t)] < \ln\varphi.
\]
\hfill $\square$

\textbf{Interpretation.}
This proves that CSL-compatible updates always exist locally provided the entropy functional varies continuously along the step path. Hence CSL enforcement by backtracking or line search is mathematically justified, not merely heuristic.

\subsection*{A.5. Operator Norm Bound for the Recursive Update}

To express recursive stability in operator form, define the state vector
\[
X(t) := \big(\{\theta_p(t)\}_{p\in\mathbb{P}_{\mathrm{active}}}, \{\phi_p(t)\}_{p\in\mathbb{P}_{\mathrm{active}}}\big)
\]
in a finite-dimensional product space $\mathcal{X}_t$ equipped with norm
\[
\|X\|_{\mathcal{X}}^2
:=
\sum_{p}\|\theta_p\|^2 + \lambda \sum_p |\phi_p|^2
\]
for some $\lambda>0$.

\textbf{Assumption A.4 (Local Lipschitz recursion).}
The recursive map $\Xi:\mathcal{X}_t\to \mathcal{X}_{t+1}$ is locally Lipschitz with constant $L_\Xi$:
\[
\|\Xi(X)-\Xi(Y)\|_{\mathcal{X}}
\le
L_\Xi \|X-Y\|_{\mathcal{X}}.
\]

\textbf{Proposition A.5 (One-step perturbation growth).}
Under Assumption A.4,
\[
\|X(t+1)-Y(t+1)\|_{\mathcal{X}}
\le
L_\Xi \|X(t)-Y(t)\|_{\mathcal{X}}.
\]

\textbf{Proof.}
Immediate from the Lipschitz assumption. \hfill $\square$

\textbf{Corollary A.5.1 (Recursive stability).}
If $L_\Xi<1$, then the recursive update is a contraction and satisfies
\[
\|X(t+n)-Y(t+n)\|_{\mathcal{X}}
\le
L_\Xi^n \|X(t)-Y(t)\|_{\mathcal{X}}.
\]

\textbf{Interpretation.}
When the payload dynamics are chosen so that the harness-level recursion is contractive, nearby cognitive states converge exponentially under iteration. This provides a rigorous form of recursive stability for the harness.

\subsection*{A.6. Norm Bound for the Hermitian EEG Interface}

The observable learning signal is defined by
\[
\mathrm{EEG}_{\mathrm{Learn}}(t)
=
\mathcal{F}^{-1}\!\left(
\langle \Psi(t)\vert \hat{H}_\Xi \vert \Psi(t)\rangle
\right),
\]
where $\hat{H}_\Xi$ is Hermitian. \cite{neuroplasticityMT,QuantumEEG1}

\textbf{Proposition A.6 (Reality of the expectation value).}
If $\hat{H}_\Xi$ is Hermitian, then
\[
\langle \Psi \vert \hat{H}_\Xi \vert \Psi \rangle \in \mathbb{R}
\]
for every $\Psi$ in the domain of $\hat{H}_\Xi$.

\textbf{Proof.}
For Hermitian $\hat{H}_\Xi$,
\[
\langle \Psi \vert \hat{H}_\Xi \vert \Psi \rangle
=
\overline{\langle \Psi \vert \hat{H}_\Xi \vert \Psi \rangle},
\]
hence the scalar equals its complex conjugate and is therefore real. \cite{HermitianOps,HermiticityNotes} \hfill $\square$

\textbf{Proposition A.7 (Expectation-value bound).}
If $\hat{H}_\Xi$ is bounded on $\mathcal{H}_t$, then
\[
\big|\langle \Psi \vert \hat{H}_\Xi \vert \Psi \rangle\big|
\le
\|\hat{H}_\Xi\|_{\mathrm{op}} \,\|\Psi\|^2.
\]

\textbf{Proof.}
By Cauchy--Schwarz,
\[
\big|\langle \Psi \vert \hat{H}_\Xi \vert \Psi \rangle\big|
\le
\|\Psi\|\,\|\hat{H}_\Xi \Psi\|
\le
\|\Psi\|\,\|\hat{H}_\Xi\|_{\mathrm{op}}\,\|\Psi\|
=
\|\hat{H}_\Xi\|_{\mathrm{op}}\,\|\Psi\|^2.
\]
\hfill $\square$

\textbf{Corollary A.7.1 (EEG amplitude bound).}
If $\mathcal{F}^{-1}$ is realized unitarily (or up to a known normalization constant $C_{\mathcal{F}}$), then
\[
\|\mathrm{EEG}_{\mathrm{Learn}}\|_{L^2}
\le
C_{\mathcal{F}}\,\|\hat{H}_\Xi\|_{\mathrm{op}}\,\|\Psi\|^2.
\]

\textbf{Interpretation.}
A bounded Hermitian interface yields a bounded observable signal. Hence signal amplitude cannot diverge unless either the state norm or the operator norm diverges.

\subsection*{A.7. Prime-Local Decomposition Bound}

Suppose the Hermitian operator admits a prime-local plus coupling decomposition
\[
\hat{H}_\Xi
=
\sum_{p\in\mathbb{P}_{\mathrm{active}}} \hat{H}_p
+
\sum_{\substack{p,q\in\mathbb{P}_{\mathrm{active}}\\ p<q}} \hat{H}_{pq},
\]
where each $\hat{H}_p$ acts locally at address $p$ and $\hat{H}_{pq}$ encodes pairwise coupling.

\textbf{Proposition A.8 (Operator norm subadditivity bound).}
If all summands are bounded, then
\[
\|\hat{H}_\Xi\|_{\mathrm{op}}
\le
\sum_p \|\hat{H}_p\|_{\mathrm{op}}
+
\sum_{p<q} \|\hat{H}_{pq}\|_{\mathrm{op}}.
\]

\textbf{Proof.}
The operator norm is subadditive:
\[
\|A+B\|_{\mathrm{op}} \le \|A\|_{\mathrm{op}} + \|B\|_{\mathrm{op}}.
\]
Applying this inductively to the decomposition gives the result. \hfill $\square$

\textbf{Corollary A.8.1.}
If
\[
\sup_p \|\hat{H}_p\|_{\mathrm{op}} \le h_1,
\qquad
\sup_{p<q}\|\hat{H}_{pq}\|_{\mathrm{op}} \le h_2,
\]
and $N=|\mathbb{P}_{\mathrm{active}}|$, then
\[
\|\hat{H}_\Xi\|_{\mathrm{op}}
\le
N h_1 + \frac{N(N-1)}{2} h_2.
\]

This gives an explicit scaling law for the observable interface as a function of harness complexity.

\subsection*{A.8. Bound on Phase-Coherence Updates}

Suppose phase updates take the form
\[
\phi_p(t+1)
=
\phi_p(t) + \alpha\, \mathcal{C}_p(t)\,\partial_{\phi_p}\mathcal{L}(\Psi(t)),
\]
where $\alpha>0$ is a step size, $\mathcal{C}_p(t)$ is a coherence factor, and $\mathcal{L}$ is a learning functional.

\textbf{Proposition A.9 (Uniform phase increment bound).}
If
\[
|\mathcal{C}_p(t)| \le C_{\max},
\qquad
\big|\partial_{\phi_p}\mathcal{L}(\Psi(t))\big| \le G_{\max}
\]
for all active $p$, then
\[
|\phi_p(t+1)-\phi_p(t)|
\le
\alpha C_{\max} G_{\max}.
\]

\textbf{Proof.}
Immediate from the update formula. \hfill $\square$

\textbf{Corollary A.9.1.}
If the EEG interface depends Lipschitz-continuously on the phases with constant $L_\phi$, then the induced signal perturbation is bounded by
\[
\|\Delta \mathrm{EEG}\|
\le
L_\phi \,\alpha C_{\max} G_{\max}\,\sqrt{N},
\]
where $N=|\mathbb{P}_{\mathrm{active}}|$.

This makes explicit how phase control and step size regulate measurement variability.

\subsection*{A.9. EchoBraid Norm Bound}

For the neurodiversity-braided state
\[
\mathrm{EchoBraid}_\nu(t)
=
\sum_{n=1}^{\infty}\psi_{p_n}^{1-\nu}(t)\otimes e^{i\theta_{p_n}^{\nu}(t)},
\]
consider the finite active truncation
\[
\mathrm{EchoBraid}_{\nu,N}(t)
=
\sum_{n=1}^{N}\psi_{p_n}^{1-\nu}(t)\otimes e^{i\theta_{p_n}^{\nu}(t)}.
\]

\textbf{Proposition A.10 (Finite truncation bound).}
If
\[
\sum_{n=1}^{\infty}\|\psi_{p_n}(t)\|^2 < \infty,
\]
then the finite truncations converge in norm, and
\[
\|\mathrm{EchoBraid}_{\nu,N}(t)\|^2
\le
\sum_{n=1}^N \|\psi_{p_n}(t)\|^{2(1-\nu)}.
\]

\textbf{Proof sketch.}
Each phase factor has unit modulus, so it does not change norms. Square-summability of the amplitude family implies Cauchy convergence of the finite partial sums in the Hilbert norm. The stated upper bound follows from triangle-type control on the truncated series. \hfill $\square$

\textbf{Interpretation.}
EchoBraid is mathematically well-defined whenever the braided amplitude family is square-summable on active addresses.

\subsection*{A.10. CSL-Compatible Operator Design Criterion}

The previous results suggest a design principle.

\textbf{Proposition A.11 (Sufficient criterion for bounded observable recursion).}
Assume:

\begin{enumerate}
  \item The recursive operator $\Xi$ is locally Lipschitz with constant $L_\Xi$.
  \item CSL enforcement guarantees
  \[
  S[\Psi(t+1)]-S[\Psi(t)] < \ln\varphi
  \]
  at each learning step.
  \item The Hermitian operator $\hat{H}_\Xi$ is bounded with
  \[
  \|\hat{H}_\Xi\|_{\mathrm{op}} \le H_{\max}.
  \]
  \item The state norm remains bounded:
  \[
  \|\Psi(t)\| \le M
  \qquad \text{for all } t.
  \]
\end{enumerate}

Then the observable EEG-like signal remains uniformly bounded:
\[
\|\mathrm{EEG}_{\mathrm{Learn}}(t)\|_{L^2}
\le
C_{\mathcal{F}} H_{\max} M^2.
\]

\textbf{Proof.}
Apply Proposition A.7 and Corollary A.7.1 at each time step. \hfill $\square$

\textbf{Interpretation.}
Stable recursion, bounded entropy change, and bounded operator norm suffice to keep the measurement interface under control.

\subsection*{A.11. Practical Proof Obligations for Future Work}

To upgrade the appendix from defensive publication to journal-grade formalism, the following proof obligations remain:

\begin{enumerate}
  \item A sharper derivation of the CSL threshold $\ln\varphi$, rather than treating it as a structurally motivated axiom. \cite{GoldenEEG}
  \item Existence and uniqueness theorems for recursive trajectories under a specified payload class $\Xi$.
  \item Spectral theorems for prime-local or prime-coupled constructions of $\hat{H}_\Xi$.
  \item Identifiability results for recovering address weights or phases from empirical EEG data.
\end{enumerate}

\subsection*{A.12. Summary of Explicit Bounds}

For convenience, we collect the core bounds here:

\[
0 \le S[\Psi(t)] \le \ln |\mathbb{P}_{\mathrm{active}}(t)|
\]

\[
|S(\pi)-S(\tilde{\pi})| \le C_N \|\pi-\tilde{\pi}\|_1
\]

\[
\|X(t+1)-Y(t+1)\|_{\mathcal{X}}
\le
L_\Xi \|X(t)-Y(t)\|_{\mathcal{X}}
\]

\[
\big|\langle \Psi \vert \hat{H}_\Xi \vert \Psi \rangle\big|
\le
\|\hat{H}_\Xi\|_{\mathrm{op}} \,\|\Psi\|^2
\]

\[
\|\hat{H}_\Xi\|_{\mathrm{op}}
\le
\sum_p \|\hat{H}_p\|_{\mathrm{op}}
+
\sum_{p<q}\|\hat{H}_{pq}\|_{\mathrm{op}}
\]

\[
|\phi_p(t+1)-\phi_p(t)|
\le
\alpha C_{\max} G_{\max}
\]

\[
\|\mathrm{EEG}_{\mathrm{Learn}}(t)\|_{L^2}
\le
C_{\mathcal{F}}\,\|\hat{H}_\Xi\|_{\mathrm{op}}\,\|\Psi(t)\|^2
\]

These are the core quantitative statements presently derived by the harness formalism.

\nocite{*}
\bibliographystyle{unsrt}  
\bibliography{references}  


\end{document}
